% ──────────────────────────────────────────────────────────────────────────────
%  Section 4 — Peak-to-peak modular spread of the phase currents
% ──────────────────────────────────────────────────────────────────────────────

\section{\texorpdfstring{Peak-to-peak modular spread $\Delta_{\max}(S,\ell)$}{Peak-to-peak modular spread Delta_max(S, ell)}}
\label{sec:spread}

The largest modular deviation between phases at the local terminal is
the peak-to-peak spread of the three magnitudes
\autoref{eq:Ik-mod}. With the line-only triplet
$\bigl(\xi_{k}(\ell),\,\eta_{k}(\ell),\,\zeta_{k}(\ell)\bigr)$ defined in
\autoref{eq:xi-eta-zeta},
%
\begin{equation}
\label{eq:Deltamax}
\boxed{\;
\begin{aligned}
\Delta_{\max}(S,\ell)
\;=\;\;& \max_{k\in\{A,B,C\}}
   \sqrt{\xi_{k}\,|S|^{2} + 2\,\Re[\eta_{k}]\,P + 2\,\Im[\eta_{k}]\,Q + \zeta_{k}} \\[2pt]
{}\;-\;& \min_{k\in\{A,B,C\}}
   \sqrt{\xi_{k}\,|S|^{2} + 2\,\Re[\eta_{k}]\,P + 2\,\Im[\eta_{k}]\,Q + \zeta_{k}}.
\end{aligned}
\;}
\end{equation}
%
\noindent The two extrema are taken over the same scalar expression
under the radical, evaluated with the per-phase constants
$(\xi_{A},\eta_{A},\zeta_{A})$, $(\xi_{B},\eta_{B},\zeta_{B})$, and
$(\xi_{C},\eta_{C},\zeta_{C})$. The whole right-hand side depends on
$\ell$ through the nine line-only scalars and on the operating point
through $(P,Q) = (\Re S,\Im S)$. No other quantity enters.

A useful corollary expresses the pairwise modulus difference in closed
form. Subtracting \autoref{eq:Ik-modsq-real} for two phases $i,j$
removes the constant of $V_{s}$ and yields
%
\begin{equation}
\label{eq:pairwise}
\bigl|I_{i}\bigr|^{2}
-
\bigl|I_{j}\bigr|^{2}
\;=\;
\bigl(\xi_{i}-\xi_{j}\bigr)\,|S|^{2}
\;+\;
2\,\Re\!\bigl[\eta_{i}-\eta_{j}\bigr]\,P
\;+\;
2\,\Im\!\bigl[\eta_{i}-\eta_{j}\bigr]\,Q
\;+\;
\bigl(\zeta_{i}-\zeta_{j}\bigr).
\end{equation}
%
\noindent which is a pure scalar polynomial of degree two in $(P,Q)$,
linear in the differences of the line constants only. The
peak-to-peak spread \autoref{eq:Deltamax} is then equivalently the
largest absolute value of $\bigl||I_{i}|-|I_{j}|\bigr|$ over the three
pairs $\{AB, BC, CA\}$.

\subsection{Two operational corollaries}

\paragraph{Light-loading asymptote.}
For small $|S|$ the quadratic term in $|S|^{2}$ is dominated by the
linear and constant parts, and \autoref{eq:Ik-mod} reduces to
%
\begin{equation}
\label{eq:lightloading}
\bigl|I_{k}(S,\ell)\bigr|
\;\xrightarrow{|S|\to 0}\;
\sqrt{\zeta_{k}(\ell)}
\;+\;
\frac{\Re[\eta_{k}(\ell)]\,P + \Im[\eta_{k}(\ell)]\,Q}{\sqrt{\zeta_{k}(\ell)}}
\;+\;
O(|S|^{2}).
\end{equation}
%
\noindent At true no load ($S=0$) the per-phase magnitude is exactly
$\sqrt{\zeta_{k}(\ell)}=|b_{k}(\ell)|$, the line-charging current of
each phase. The peak-to-peak spread then collapses to
$\max_{k}|b_{k}(\ell)|-\min_{k}|b_{k}(\ell)|$, which depends on the line
only.

\paragraph{Heavy-loading asymptote.}
For $|S|$ large the quadratic term in $|S|^{2}$ dominates and
%
\begin{equation}
\label{eq:heavyloading}
\bigl|I_{k}(S,\ell)\bigr|
\;\xrightarrow{|S|\to\infty}\;
\sqrt{\xi_{k}(\ell)}\,|S|
\;+\;
\frac{\Re[\eta_{k}(\ell)]\,P + \Im[\eta_{k}(\ell)]\,Q}{\sqrt{\xi_{k}(\ell)}\,|S|}
\;+\;
O\!\bigl(|S|^{-3}\bigr).
\end{equation}
%
\noindent The relative spread, that is the absolute spread divided by
the average magnitude, then approaches the spread of
$\sqrt{\xi_{k}(\ell)}$ across the three phases. Because
$|a_{k}|=|p_{k}|/(V_{s}|\alpha|)$ is set by the moduli of the
phase-projected $\matB^{-1}\vu$, the heavy-load relative unbalance is
controlled purely by the asymmetry of $\vp(\ell)$ across the three
phases.
